% Part: normal-modal-logic
% Chapter: axioms-systems
% Section: logics-proofs

\documentclass[../../../include/open-logic-section]{subfiles}

\begin{document}

\olfileid{nml}{prf}{prf}

\olsection{\usetoken{P}{derivation} و دستگاه‌های موجهاتی}

نخست تعریف می‌کنیم !!a{derivation} برای منطق‌های موجهات عادی چیست.
به‌طور تقریبی، !!a{derivation} دنباله‌ای از !!{formula}هاست که هر
!!{element} آن یا (نمونهٔ جانشینیِ) یکی از چند \emph{اصل} است یا با یکی
از چند قاعدهٔ استنتاج از !!{element}های پیشین نتیجه می‌شود. برای
منطق‌های موجهات عادی، همهٔ نمونه‌های همان‌گویی‌ها
\iftag{prvDiamond}{، \Ax{K} و~\Dual{}}{ و~\Ax{K}} اصل به شمار می‌آیند.
حاصل دستگاه موجهاتی~$\Log{K}$، یعنی کوچک‌ترین منطق موجهات عادی، است.
بااین‌حال، ممکن است بخواهیم برای به‌دست‌آوردن دستگاه‌های دیگر اصل‌های
دیگری بیفزاییم. قواعد همیشه وضع مقدم~\MP{} و ضرورت‌بخشی~\Nec{} هستند.

\begin{defn}
  با داشتن دستگاه موجهاتی~$\Log{K} !A_1 \dots !A_n$ و
  !!a{formula}~$!B$، می‌گوییم~$!B$ در~$\Log{K} !A_1 \dots !A_n$
  \emph{!!{derivable}} است، به نماد
  $\Log{K} !A_1 \dots !A_n \Proves !B$، اگر و تنها اگر
  !!{formula}های~$!C_1$، \dots، $!C_k$ وجود داشته باشند چنان‌که
  $!C_k = !B$ و هر~$!C_i$ یا یک نمونهٔ همان‌گویی، یا نمونه‌ای از یکی
  از~$\Ax{K}$،\iftag{prvDiamond}{ $\Dual$،}{} $!A_1$، \dots، $!A_n$
  باشد، یا با قواعد~\MP{} یا~\Nec{} از !!{formula}های پیشین نتیجه شود.
\end{defn}

گزارهٔ زیر به ما امکان می‌دهد با ارائهٔ یک~$\Sigma$-!!{derivation} از
$!B$ نشان دهیم~$!B \in \Sigma$.

\begin{prop}
  $\Log{K} !A_1 \dots !A_n = \Setabs{!B}{\Log{K} !A_1 \dots !A_n \Proves !B}$.
\end{prop}

\begin{proof}
  با استقرا بر طول !!{derivation}ها نشان می‌دهیم
  $\Setabs{!B}{\Log{K} !A_1 \dots !A_n \Proves !B} \subseteq
  \Log{K} !A_1 \dots !A_n$.

  اگر !!{derivation}~$!B$ طول~$1$ داشته باشد، تنها شامل یک
  !!{formula} است. آن !!{formula} نمی‌تواند با~\MP{} یا~\Nec{} از
  فرمول‌های پیشین نتیجه شده باشد؛ پس باید یک نمونهٔ همان‌گویی، نمونه‌ای
  از~\Ax{K}،\iftag{prvDiamond}{ $\Dual$،}{} یا نمونه‌ای از~$!A_1$،
  \dots،~$!A_n$ باشد. اما~$\Log{K}!A_1\dots!A_n$ این‌ها را نیز شامل
  می‌شود؛ پس~$!B \in \Log{K}!A_1 \dots !A_n$.

  اگر !!{derivation}~$!B$ طولی بزرگ‌تر از~$1$ داشته باشد، افزون بر این
  ممکن است~$!B$ با~\MP{} یا~\Nec{} از !!{formula}هایی به دست آمده باشد
  که خط آخر !!{derivation} نیستند. اگر~$!B$ با~\MP{} از~$!C$ و
  $!C \lif !B$ نتیجه شود، بنا بر فرض استقرا~$!C$ و
  $!C \lif !B \in \Log{K}!A_1 \dots !A_n$. اما هر منطق موجهات تحت وضع
  مقدم بسته است، پس~$!B \in \Log{K}!A_1 \dots !A_n$. اگر
  $!B \ident \Box !C$ با~\Nec{} از~$!C$ نتیجه شود، بنا بر فرض استقرا
  $!C \in \Log{K}!A_1 \dots !A_n$. اما هر منطق موجهات عادی تحت~$\Nec$
  بسته است؛ پس~$!B \in \Log{K}!A_1\dots!A_n$.

  شمول عکس با نشان‌دادن این امر به دست می‌آید که
  $\Sigma = \Setabs{!B}{\Log{K} !A_1 \dots !A_n \Proves !B }$ یک منطق
  موجهات عادی شامل همهٔ نمونه‌های~$!A_1$، \dots، $!A_n$ است و با توجه
  به اینکه~$\Log{K} !A_1 \dots !A_n$ بنا به تعریف کوچک‌ترین منطق از
  این نوع است.
  \begin{enumerate}
    \item هر همان‌گویی~$!B$ یک نمونهٔ همان‌گویی است؛ پس
      $\Log{K}!A_1\dots!A_n \Proves !B$، و~$\Sigma$ همهٔ همان‌گویی‌ها
      را شامل می‌شود.
    \item اگر~$\Log{K}!A_1\dots!A_n \Proves !C$ و
      $\Log{K}!A_1\dots!A_n \Proves !C \lif !B$، آنگاه
      $\Log{K}!A_1\dots!A_n \Proves !B$: !!{derivation}~$!C$ را با
      اثباتِ~$!C \lif !B$ ترکیب کنید و خط~$!B$ را بیفزایید.
      خط آخر با~\MP{} توجیه می‌شود. پس~$\Sigma$ تحت وضع مقدم بسته است.
    \item اگر~$!B$ !!a{derivation} داشته باشد، هر نمونهٔ جانشینی~$!B$
      نیز !!a{derivation} دارد: جانشینی را بر هر !!{formula} در
      !!{derivation} اعمال کنید. (تمرین: با استقرا بر طول
      !!{derivation}ها ثابت کنید نتیجه نیز !!{derivation} درستی است.)
      پس~$\Sigma$ تحت جانشینی یکنواخت بسته است. (اکنون ثابت کرده‌ایم
      $\Sigma$ همهٔ شرایط یک منطق موجهات را برآورده می‌کند.)
    \item داریم~$\Log{K}!A_1\dots!A_n \Proves \Ax{K}$، پس
      $\Ax{K} \in \Sigma$.
    \tagitem{prvDiamond}{داریم
      $\Log{K}!A_1\dots!A_n \Proves \Dual$، پس~$\Dual \in \Sigma$.}{}
    \item اگر~$\Log{K}!A_1\dots!A_n \Proves !C$، خط اضافی~$\Box !C$
      با~\Nec{} توجیه می‌شود. بنابراین~$\Sigma$ تحت~\Nec{} بسته است؛
      پس~$\Sigma$ عادی است.
  \end{enumerate}
\end{proof}

\end{document}
